\subsection{Why the generic marking cannot simply be omitted}
\label{sec:archive-unmarked-curve}

\begin{proposition}[F149-UNMARKED-CURVE]
\label{prop:archive-unmarked-curve}
Even the full relative curve diagram
\[
 I\ \subset\ N\ \subset\ \Pi\ \longrightarrow\ G
       \longrightarrow G_{\mathbb Q},                       \tag{UC.1}
\]
with its identified arithmetic quotient, marked cusp inertia and
cyclotomic action, need not determine the geometric tangent-section
subset invariantly under its profinite automorphisms. This remains
false if one retains a reference constant tangent section and a
rational base section of $G\to G_{\mathbb Q}$.

More precisely, take
\[
 S=\mathbb G_{m,\mathbb Q}=\operatorname{Spec}\mathbb Q[t,t^{-1}],
 \quad C=\mathbb P^1_{\mathbb Q}\setminus\{0,1,\infty\},
 \quad U=S\times_{\mathbb Q}C,
\]
with marked relative cusp $0$ and its constant normal basis $b$.
There is an automorphism of (UC.1), acting identically on $I$ and
$G_{\mathbb Q}$, which carries the section class associated with the
nowhere-zero normal vector $t b$ to a class not associated with
any global nowhere-zero normal vector.
The automorphism is allowed to act nontrivially on $G$; no assertion
is made with $G$ fixed pointwise, with its scheme realization fixed,
or with the generic Galois marking retained.
\end{proposition}

\begin{proof}
Write $\Gamma=G_{\mathbb Q}$. Use the rational point $t=1$ to
split the arithmetic exact sequence for $S$. The geometric
fundamental group of $S$ is
\[
 A=\pi_1(S_{\overline{\mathbb Q}})=\widehat{\mathbb Z}(1),
 \qquad G=A\rtimes\Gamma.                                  \tag{UC.2}
\]
The first statement follows from the punctured-curve presentation
in \cite[Expos\'e XIII, Corollaire~2.12, pp.~290--292]{SGA1};
the compatible covers $v^n=t$ identify the Galois action on $A$
with the cyclotomic action. They have the compatible points $v=1$
over $t=1$. These covers and their Kummer normalization are
\cite[Expos\'e XI, Proposition~6.1 and \S6.2--6.5,
pp.~233--235]{SGA1}.

The constant tangent basis $b$ gives a section of the cusp
normalizer of $\pi_1(C)\to\Gamma$. Choose the resulting tangential
basepoint and write
\[
 \pi_1(C)=\Delta\rtimes\Gamma,
 \quad \Delta=\pi_1(C_{\overline{\mathbb Q}}),
 \quad I\subset\Delta.
\]
The normalizer at that cusp is $I\rtimes\Gamma$.
Here the normal-line comparison and its section are exactly
F149-RELATIVE, applied over $\operatorname{Spec}\mathbb Q$;
geometric self-normalization is also
\cite[Proposition~16, p.~8]{StixCuspidal2012}.
The action of $\Gamma$ on $I=\widehat{\mathbb Z}(1)$ is again
cyclotomic.

For completeness, the full product group, rather than merely
its normalizer, is available. The natural map
\[
 \pi_1(U)\longrightarrow
        \pi_1(S)\mathbin{\times}_{\Gamma}\pi_1(C)             \tag{UC.3}
\]
is an isomorphism. On the geometric kernels it is the
K\"unneth isomorphism of
\cite[Expos\'e XIII, Proposition~4.6, pp.~310--311]{SGA1}.
Its hypotheses hold for the two smooth curves over
$\overline{\mathbb Q}$: in characteristic zero all primes are
allowed, and resolution in dimension at most one suffices.
The arithmetic fundamental exact sequence
\cite[Expos\'e IX, Th\'eor\`eme~6.1, p.~195]{SGA1} then makes
(UC.3) a homomorphism of short exact sequences which is an
isomorphism on the kernel and on the quotient $\Gamma$.
It is consequently an isomorphism of profinite groups.
With the chosen sections this gives
\[
 \Pi=(A\times\Delta)\rtimes\Gamma,
 \qquad N=(A\times I)\rtimes\Gamma.                         \tag{UC.4}
\]
Indeed, $A$ centralizes $\Delta$, and within $\Delta$ the
normalizer of $I$ is $I$; the chosen arithmetic section normalizes
$I$. Thus the second group in (UC.4) is the actual full
normalizer in $\Pi$. Projection to $G$ forgets the $I$ coordinate.
It also identifies $N$ with $I\rtimes G$, where $G$ acts on $I$
through its cyclotomic quotient $\Gamma$.

Choose the profinite unit
\[
 u\in\widehat{\mathbb Z}^{\times},\qquad
 u_2=3\in\mathbb Z_2^\times,\quad
 u_\ell=1\in\mathbb Z_\ell^\times\quad(\ell\ne2).           \tag{UC.5}
\]
Multiplication by $u$ on $A$ commutes with its cyclotomic
$\Gamma$-action. Therefore
\[
 \begin{split}
 \alpha_G(a,g)&=(ua,g),\\
 \alpha_\Pi(a,\delta,g)&=(ua,\delta,g)
 \end{split}                                               \tag{UC.6}
\]
are continuous group automorphisms, with inverse obtained by
replacing $u$ by $u^{-1}$. They commute with the projections in
(UC.1), fix its arithmetic quotient, and preserve $N$.
The restriction $\alpha_N$ fixes $I$ pointwise. They also fix
the chosen $\Gamma$ section, the constant tangent section, and
indeed all geometric cusp-inertia subgroups lying in $\Delta$.

Identify $A$ with $I$ through their common Kummer normalization.
Write an element of $N=I\rtimes G$ as $(c;(a,g))$.
Relative to the constant tangent section, the normal vector
$t b$ has cocycle
\[
 c_t(a,g)=a.                                               \tag{UC.7}
\]
This can be checked at every finite level using $v^n=t$:
the geometric inertia acts by its tautological $\mu_n$ quotient,
and the section $t=1$, lifted by $v=1$, has zero cocycle.
Equivalently, (UC.7) satisfies the cocycle identity directly,
since the two inertia modules have the same action of $g$.
Consequently the corresponding section is
$s_t(a,g)=(a;(a,g))$.

Transporting a section by an automorphism of the diagram means
$s\mapsto\alpha_N\circ s\circ\alpha_G^{-1}$. Thus the
transported $s_t$ is represented by
\[
 c'_t(a,g)=u^{-1}a.                                       \tag{UC.8}
\]
Every global nowhere-zero normal vector is of the form
$f b$, where
\[
 f\in\mathbb Q[t,t^{-1}]^\times=\mathbb Q^\times t^{\mathbb Z}.
                                                               \tag{UC.9}
\]
To see the last equality directly, the difference between the
highest and lowest exponent of a nonzero Laurent polynomial
is additive under multiplication over a field. A unit and its
inverse therefore both have exponent difference zero.
Write $f=a_0t^m$ with $a_0\in\mathbb Q^\times$, $m\in\mathbb Z$.
Kummer functoriality shows that its section cocycle restricts
to $A$ as multiplication by $m$; the constant factor has
zero restriction to $A$.

Restriction to $A$ is unchanged by $I$-conjugacy of sections:
$I$ is abelian and $A$ acts trivially on it, so every coboundary
restricts to zero. If (UC.8) were geometrically represented,
we would therefore have $m=u^{-1}$ in
$\widehat{\mathbb Z}$ for some ordinary integer $m$.
Its $2$-adic component would imply $3m=1$ in $\mathbb Z_2$,
hence $3m=1$ in $\mathbb Z$, which is impossible.
This proves that the transported section class is nongeometric.
A uniform reversal of normal/conormal inertia conventions changes
both slopes by the same sign and does not change the contradiction.
\end{proof}

\begin{corollary}[F149-NO-FUNCTORIAL-UNMARKING]
\label{cor:archive-no-unmarking}
On the class of the realizable diagrams in (UC.1), the full
generic arithmetic marking required in
F149-MARKED-RECOGNITION cannot be recovered functorially in a way
that both recovers its geometric realization and lifts every
isomorphism of the smaller diagram. In particular no
isomorphism-invariant recognition rule on the smaller data can
select all and only geometric tangent sections.
\end{corollary}

\begin{proof}
Such a recovery would lift the automorphism (UC.6) to an
isomorphism of full marked generic diagrams. By
F149-MARKED-RECOGNITION that isomorphism preserves the geometric
section subset. Proposition~\ref{prop:archive-unmarked-curve}
shows that its induced automorphism of the smaller diagram does
not preserve that subset, a contradiction.
The assertion concerns functorial recovery from the specified
profinite data. It does not exclude adding the scheme $S$, its
rational-point markings, or the generic Galois group as further
input, nor a noncanonical choice that is not respected by
all diagram isomorphisms.
\end{proof}
